\subsection{The monodromy condition and the Fontaine comparison}

\begin{proposition}[B-RD4: Lead B' needs an additional hypothesis]\label{prop:B-RD4}
\status{Refuted as stated.}
For every odd prime $p$, there is a height-two Barsotti--Tate group over
$K=\mathbf Q_p(\zeta_p)$ satisfying the manuscript's monodromy condition
whose rational Tate module is not potentially Hodge--Tate, and therefore
is not potentially semistable. Its manuscript inertia character is the
constant function $2$. Thus the monodromy condition in 9:11 cannot be
identified with potential semistability, and Fontaine's usual
Weil--Deligne representation need not be defined for its Tate module.
\end{proposition}
\begin{proof}
Write $\chi:G_K\to\mathbf Z_p^\times$ for the cyclotomic character.
Since $\mu_p\subset K$, it takes its values in $1+p\mathbf Z_p$. Put
\[
 \rho(g)=\begin{pmatrix}1&p^{-1}\log\chi(g)\\0&1\end{pmatrix}.
\]
This is a continuous representation on $T=\mathbf Z_p^2$, because the
logarithm is additive on $1+p\mathbf Z_p$ and takes values in
$p\mathbf Z_p$. Over a characteristic-zero field, the equivalence
between finite etale group schemes and finite continuous Galois modules
applied to $T/p^nT$ gives compatible finite etale group schemes $M[p^n]$.
The transition maps are induced by multiplication by $p$ on the lattices,
and the resulting inductive system is a height-two Barsotti--Tate group.
The invariant first coordinate and the quotient both give the constant
group $\mathbf Q_p/\mathbf Z_p$. They extend over $\mathcal O_K$, so the
specified two-step filtration satisfies the manuscript condition and
its character is $1+1=2$.

For $V=T\otimes\mathbf Q_p$, the Sen operator is
\[
 \Theta_V=\frac1p\begin{pmatrix}0&1\\0&0\end{pmatrix}.
\]
Indeed the logarithm of a unipotent matrix $1+aN$, with $N^2=0$, is
$aN$, and division by $\log\chi(g)$ gives the displayed operator.
It remains nonzero and nilpotent after restriction to any finite
extension. Sen's criterion says that a representation is Hodge--Tate
exactly when this operator is diagonalizable with integral eigenvalues;
the identical unscaled example and the criterion are explicitly given
in \cite[\S II.1.2, p.~9 of the author's version]{Berger2004Introduction}.
The operator is not diagonalizable. Potentially semistable
representations are potentially Hodge--Tate, so no finite extension
makes this representation semistable.
\end{proof}

\begin{proposition}[The corrected Fontaine dictionary]
\label{prop:B-RD4-dictionary}
\status{Proved with the additional potential semistability hypothesis.}
Let $M$ satisfy the manuscript monodromy condition, and assume in
addition that $V=T_p(M)\otimes_{\mathbf Z_p}\mathbf Q_p$ is potentially
semistable. For arbitrary perfect residue field, its manuscript
character is the inertia trace of Fontaine's $p$-module of Deligne
$D_{\mathrm{pst}}(V)$. When the residue field is finite, this is the
inertia character of the usual Weil--Deligne representation. With
\[
 D_{\mathrm{cov}}(G)=D_{\mathrm{contra}}(G^D),
\]
the precise comparison for a good graded group $G$ over a finite
extension $L/K$ is
\[
 D_{\mathrm{cov}}(G_{k_L})[1/p]
       \simeq D_{\cris,L}(T_pG\otimes\mathbf Q_p)(-1).
                                                        \tag{B-RD4.1}
\]
The Tate twist in this formula does not change the inertia trace.
\end{proposition}
\begin{proof}
Choose a finite Galois extension $L/K$ over which $V$ is semistable
and every graded group has good reduction. All subrepresentations
and quotients in the filtration are semistable over $L$.
Crystalline comparison for a BT group, with the indicated
covariant-versus-contravariant convention, gives (B-RD4.1)
functorially and compatibly with descent; see
\cite[Lemma~2.2.4 and Corollary~2.2.6, author pp.26--27; Appendix~A.2
and~A.6]{Kisin2006Crystalline}. One can check the convention on
$\mathbf Q_p/\mathbf Z_p$ and $\mu_{p^\infty}$: their covariant
Frobenius slopes are respectively one and zero, whereas the
crystalline Tate-module slopes are zero and minus one.

The semistable period functor is exact on semistable
representations. Thus its filtration has graded terms the
crystalline period modules of the good graded groups, and traces
of every finite descent-inertia element add. Exactness and the
passage from the completed maximal unramified extension back to
$K$ are treated in \cite[\S\S5.1.5--5.1.7,
pp.156--158]{Fontaine1994Semi}; alternatively the exact
$D_{\mathrm{pst}}$ functor is stated in
\cite[\S2.3.4, p.337]{Fontaine1994WD}.
The crystalline Tate twist is unramified in this period-module
sense, so the right side of (B-RD4.1) has the same inertia trace as
its untwisted version. Summing proves the assertion.

Over an arbitrary perfect residue field the natural object is the
$p$-module of Deligne, with semilinear Galois action, Frobenius and
monodromy; its definition is
\cite[\S\S1.1.1--1.1.2, pp.322--323]{Fontaine1994WD}.
If the residue field is finite, the corresponding linear
Weil--Deligne action is
$\rho_0(w)=\varphi^{-a(w)}\rho(w)$, as in
\cite[\S1.3.5, p.329; \S2.3.7, pp.339--340]{Fontaine1994WD}.
For inertia $a(w)=0$, so the trace just computed is exactly the
Weil--Deligne inertia trace. No finite-residue-field hypothesis is
needed for the preceding Deligne-module statement.
\end{proof}

\begin{proposition}[B-RD5: the height-one case of B.4]\label{prop:B-RD5}
\status{Proved.}
Suppose $\mu_p\subset K$. For height $h=1$ and every prime $p$, the manuscript
character on prime-to-$p$ inertia elements is determined by $M(0)$.
The least index in its indexing is therefore $\nu=0$.
\end{proposition}
\begin{proof}
For $p=2$, the finite special inertia action is contained in
$(\mathbf Q_2^\times)_{\mathrm{tors}}=\{1,-1\}$. Every odd-order
inertia element therefore acts trivially. Its requested character is
identically one, so $\nu=0$ suffices. We now assume $p\geq3$.

A height-one group satisfying the monodromy condition itself acquires
good reduction. Its good model has dimension zero or one. In dimension
zero it is etale, and in dimension one its Cartier dual is etale.
Over the completed maximal unramified extension of a good-reduction
field its model is respectively $\mathbf Q_p/\mathbf Z_p$ or
$\mu_{p^\infty}$. Its finite descent inertia character therefore takes
values in the torsion subgroup of $\mathbf Z_p^\times$, namely
$\mu_{p-1}$. In dimension zero, reduction modulo $p$ gives precisely
this character. In dimension one, the Tate module acquires additionally
the cyclotomic factor, whose reduction is trivial under $\mu_p\subset K$.
Thus in both cases the finite character is the unique Teichmuller lift
of the one-dimensional mod-$p$ inertia character. The same conclusion
holds after summing characters of height-one potentially good graded
pieces: the trace on a prime-to-$p$ element is the Brauer character of
the mod-$p$ representation, and is independent of the chosen filtration.
\end{proof}
